\scanpage{56}
If
\[
\alpha:\Delta_{r'}\longrightarrow\Delta_r,\qquad \text{\uuline{$\alpha(\unsure{0|\omega})=0$}},
\]
\[
(29)\qquad\left\{\begin{aligned}
&\text{then}\qquad \alpha^*:X_{\Delta_r}\mathbin{\unsure{-|\longrightarrow}}X_{\Delta_{r'}}\qquad(\text{étale})\\
&\text{satisfies}\\
&\qquad\alpha^*(X_{\Delta_{r'}}^*)=X_{\Delta_r}^*,\quad\text{ie. }\alpha^*(\dot X_{\Delta_{r'}})=\dot X_{\Delta_r}.
\end{aligned}\right.
\]

\deleted{\unsure{In the case}}

\uline{NB} ``Morally'', $X_{\Delta_r}^*$ represents a sort of ``non-\unsure{singular|singular}'' part of $X_{\Delta_r}$, $\dot X_{\Delta_r}$ the singular part or ``boundary''. One should note that the $X_{\Delta_r}^*$ \uline{form} \uline{no} semi-simplicial space.

Nevertheless, our purpose is, to a certain extent, to reconstruct $X$ from the ``open strata'' $X_i^*$ in $X$, in the spirit of number 1 (where there are two strata $X_0=Y$ and $X_1=X$, \deleted{$\unsure{Y\subset X}$ \unsure{closed}} \inserted{($Y\subset X$ inclusion} closed), and we therefore have $Y=Y^*$, $X^*=X\setminus Y$ as in the notation of the cited passage, where there was occasion to introduce $\mathcal{V}_{\unsure{Y},X}$.
\margin{55}
